% Flussdiagramm zum Berechnungsablauf (Kapitel 3, Methodik)
% Einbindung z. B. mit \subimport{TeX/}{Flussdiagramm_methodik.tex}
\begin{tikzpicture}[scale=0.9,
    node distance=0.7cm and 0.5cm,
    font=\small,
    startstop/.style={rectangle, rounded corners, minimum width=3.2cm, minimum height=0.9cm,
        text width=3.0cm, align=center, draw=black, fill=gray!15},
    io/.style={trapezium, trapezium left angle=70, trapezium right angle=110,
        minimum width=4.1cm, minimum height=0.9cm, text width=3.8cm, align=center,
        draw=black, fill=blue!8},
    process/.style={rectangle, minimum width=3.0cm, minimum height=1.0cm,
        text width=2.9cm, align=center, draw=black, fill=orange!10},
    decision/.style={diamond, aspect=2.4, inner sep=1pt, text width=2.7cm, align=center,
        draw=black, fill=yellow!15},
    arr/.style={-Latex, thick},
]

% --- Schritt 1: Eingangsdaten ---
\node[io] (input) {Schritt 1\\Eingangsdaten:\\Wetterdaten und Gebäuderandbedingungen};

% --- Schritt 2: Betriebszustand je Stunde ---
\node[decision, below=of input] (betrieb) {Schritt 2\\Stunde innerhalb der Betriebszeit?};

\node[process, right=2.2cm of betrieb] (null) {alle Energieströme dieser Stunde = 0};

% --- Schritte 3-6: parallele Berechnung der Energieströme ---
\node[process, below=1.6cm of betrieb, xshift=-4.9cm] (heizen) {Schritt 3\\Heizleistung};
\node[process, right=of heizen] (kuehlen) {Schritt 4\\Kühlleistung\\(falls Kühl-Solltemperatur hinterlegt)};
\node[process, right=of kuehlen] (befeuchten) {Schritt 5\\Befeuchtungs-\\leistung (falls Feuchte-Sollwert hinterlegt)};
\node[process, right=of befeuchten] (ventilator) {Schritt 6\\Ventilator-\\leistung (falls Druckerhöhung und Wirkungsgrad hinterlegt)};

% --- Schritt 7: Aggregation ---
\node[process, below=1.4cm of befeuchten, xshift=-1.3cm] (aggregation) {Schritt 7\\Jahresenergiemenge und Volllaststunden je Energieart};

% --- Schritt 8: Jahresdauerlinie ---
\node[process, below=of aggregation] (jdl) {Schritt 8\\Jahresdauerlinie};


\node[startstop, below=of jdl] (ende) {Ergebnisse: Jahresenergiemengen, Volllaststunden, Jahresdauerlinien};

% --- Hilfskoordinaten fuer die Verteilerleiste (Schritte 3-6) ---
\coordinate (busJa) at ($(betrieb.south)+(0,-0.5)$);

% --- Pfeile: Eingangsdaten -> Entscheidung ---
\draw[arr] (input) -- (betrieb);

% --- nein-Zweig: keine Berechnung fuer diese Stunde ---
\draw[arr] (betrieb) -- node[above]{nein} (null);
\draw[arr] (null.east) -- ++(1.0,0) |- ($(aggregation.east)+(0.9,0)$) -- (aggregation.east);

% --- ja-Zweig: Verteilerleiste zu den Schritten 3-6 ---
\draw (betrieb.south) -- (busJa) node[midway, right]{ja};
\draw (busJa -| heizen.north) -- (busJa -| ventilator.north);
\draw[arr] (busJa -| heizen.north) -- (heizen.north);
\draw[arr] (busJa -| kuehlen.north) -- (kuehlen.north);
\draw[arr] (busJa -| befeuchten.north) -- (befeuchten.north);
\draw[arr] (busJa -| ventilator.north) -- (ventilator.north);

% --- Schritte 3-6 -> Aggregation ---
\draw[arr] (heizen.south) |- (aggregation.west);
\draw[arr] (kuehlen.south) -- (aggregation.north);
\draw[arr] (befeuchten.south) -- (aggregation.north);
\draw[arr] (ventilator.south) |- (aggregation.east);

% --- weiterer Ablauf ---
\draw[arr] (aggregation) -- (jdl);
\draw[arr] (jdl) -- (ende);


\end{tikzpicture}
